% 03_modelo.tex — Slide 3: Modelo matematico ILP
\begin{frame}{Modelo matemático (ILP)}

\begin{columns}[T]
\begin{column}{0.55\textwidth}

\textbf{Variable:} $s_t \in \{0,1\}$ para cada tripleta $t \in T$

\bigskip
\begin{block}{Formulacion ILP}
\begin{align*}
\max \quad & \sum_{t \in T} s_t \\[4pt]
\text{s.a.} \quad
  & \sum_{t \ni i} s_t \le 1 \quad \forall\, i \in X \\
  & \sum_{t \ni j} s_t \le 1 \quad \forall\, j \in Y \\
  & \sum_{t \ni k} s_t \le 1 \quad \forall\, k \in Z \\
  & s_t \in \{0,1\}
\end{align*}
\end{block}

\end{column}

\begin{column}{0.42\textwidth}

\textbf{Caracteristicas:}
\begin{itemize}
  \item $m$ variables binarias
  \item $3n$ restricciones lineales
  \item Relajacion LP \alert{no} integra en general
  \item Decisión: $\mathrm{OPT}(I) = n$?
\end{itemize}

\bigskip
\textbf{Consecuencia:}

La relajacion LP admite gap $\Rightarrow$ no basta redondear; se necesita \textbf{backtracking exacto}.

\end{column}
\end{columns}

\end{frame}
